% fixture: koma/koma-footnotes
% tier: T3
% engine: document
% covers: scrartcl's defaults (fontsize=11pt, paper=a4, typearea), KOMA-Script's footnotes (\deffootnote[1em]{1.5em}{1em}: \leftskip 1.5em, the mark in \hbox to\z@ after \hskip-.5em; \textsuperscript in the font around it, italic in \emph), \footnoterule (\kern-3pt, \vskip\z@\@plus.05fil under \raggedbottom, \kern2.6pt)
% issue: #360
\documentclass{scrartcl}
\begin{document}
Typesetting is the composition of text by means of arranging physical type or its digital equivalents. Stored letters and other symbols are retrieved and ordered according to the orthography of a language for visual display. The breaking of paragraphs into lines is a small part of this craft, but it decides how evenly the words are spread and how easily the eye moves from one line to the next.

Second paragraph with some more words to see indentation. Then a footnote\footnote{A note.} and more text, \emph{and in italics\footnote{Its mark in italics.} too}.
\end{document}
